% Template Artigo SBC (Sociedade Brasileira de Computação)
% Padrão para SBES, SBSI, SBBD, SBRC, SIBGRAPI, JBCS, etc.
% Compilação: pdflatex + bibtex + pdflatex + pdflatex
% Documentação: https://www.sbc.org.br/documentos-da-sbc/category/169-templates-para-publicacoes

\documentclass[12pt]{article}
\usepackage{sbc-template}            % classe oficial SBC (baixe em www.sbc.org.br)
\usepackage{graphicx,url}
\usepackage[brazil]{babel}
\usepackage[utf8]{inputenc}
\usepackage[T1]{fontenc}

\sloppy

\title{Título do Artigo: Subtítulo se Houver}

\author{Autor Um\inst{1}, Autor Dois\inst{2}, Autor Três\inst{1}}

\address{
    Universidade Federal do Amazonas (UFAM)\\
    Av. General Rodrigo Otávio Jordão Ramos, 6200 -- 69077-000 -- Manaus -- AM
    \nextinstitute
    Outra Instituição -- Cidade -- UF
    \email{\{autor1,autor3\}@ufam.edu.br, autor2@outra.br}
}

\begin{document}
\maketitle

\begin{abstract}
This paper presents [objective]. The methodology adopted was [...].
Results indicate that [main finding]. The contribution to the field
includes [...]. We conclude that [...].
\end{abstract}

\begin{resumo}
Este artigo apresenta [objetivo]. A metodologia adotada foi [...].
Os resultados indicam que [achado principal]. A contribuição para
o campo inclui [...]. Conclui-se que [...].
\end{resumo}

\section{Introdução}

[Contextualização. Citar trabalhos: \cite{gil2017}.]

\subsection{Problema de Pesquisa}
[Pergunta de pesquisa.]

\subsection{Objetivos}
\textbf{Geral}: [...]

\textbf{Específicos}:
\begin{itemize}
    \item [obj1];
    \item [obj2];
    \item [obj3].
\end{itemize}

\subsection{Estrutura do Artigo}
[O artigo está organizado em N seções...]

\section{Trabalhos Relacionados}

[Revisão de literatura organizada por temas. Citar trabalhos
da SBC: \cite{exemplo2022}, JBCS: \cite{outro2021}, e
internacionais: \cite{international2020}.]

\section{Metodologia}

\subsection{Tipo de Pesquisa}
A pesquisa caracteriza-se como [...] de natureza [...] com
abordagem [qualitativa/quantitativa/mista], conforme \cite{gil2017}.

\subsection{Procedimentos de Coleta}
[Descrição.]

\subsection{Procedimentos de Análise}
[Descrição.]

\section{Resultados}

\subsection{Achado 1}
[Apresentação.]

\subsection{Achado 2}
[Apresentação.]

% Exemplo de figura
\begin{figure}[ht]
\centering
\includegraphics[width=.7\textwidth]{caminho/para/figura.png}
\caption{Descrição da figura}
\label{fig:exemplo}
\end{figure}

% Exemplo de tabela
\begin{table}[ht]
\centering
\caption{Descrição da tabela}
\label{tab:exemplo}
\begin{tabular}{|l|c|r|}
\hline
\textbf{Coluna 1} & \textbf{Coluna 2} & \textbf{Coluna 3} \\
\hline
Linha 1 & Dado & Dado \\
Linha 2 & Dado & Dado \\
\hline
\end{tabular}
\end{table}

\section{Discussão}

[Diálogo dos resultados com a literatura. Implicações.]

\section{Considerações Finais}

[Síntese. Contribuições. Limitações. Agenda.]

\bibliographystyle{sbc}
\bibliography{referencias}

\end{document}
